\chapter{2010年四川卷（文）}

\section{单选题}

\begin{problem}
设集合 \(A = \{3, 5, 6, 8\}\)，集合 \(B = \{4, 5, 7, 8\}\)，则 \(A \cap B\) 等于（\quad）

\choices
    {\( \{3, 4, 5, 6, 7, 8\} \)}
    {\(\{3,6\}\)}
    {\( \{4, 7\} \)}
    {\( \{5, 8\} \)}
\end{problem}

\begin{answer}
D
\end{answer}

\begin{solution}
交集中的元素必须同时属于集合 \(A\) 和集合 \(B\). 逐个比较两集合中的元素，只有 \(5\) 和 \(8\) 同时属于两集合，所以 \(A\cap B=\{5,8\}\)，故选 D.
\end{solution}

\begin{problem}
函数 \( y = \log_2 x \) 的图象大致是（\quad）

\choices
    {\choicebitmap[width=0.15\paperwidth]{img/2010/sichuan_liberal/q2_optA.png}}
    {\choicebitmap[width=0.15\paperwidth]{img/2010/sichuan_liberal/q2_optB.png}}
    {\choicebitmap[width=0.15\paperwidth]{img/2010/sichuan_liberal/q2_optC.png}}
    {\choicebitmap[width=0.15\paperwidth]{img/2010/sichuan_liberal/q2_optD.png}}
\end{problem}

\begin{answer}
C
\end{answer}

\begin{solution}
函数 \(y=\log_2x\) 的定义域为 \(x>0\)，且底数 \(2>1\)，所以函数在定义域上单调递增. 又因为 \(\log_2 1=0\)，图象经过点 \((1,0)\)；当 \(x\to0^+\) 时，\(y\to-\infty\)，所以其左侧竖直渐近线为 \(x=0\). 具有这些特征的图象是 C，故选 C.
\end{solution}

\begin{problem}
抛物线 \( y^{2}=8x \) 的焦点到准线的距离是（\quad）

\choices
    {\(1\)}
    {\(2\)}
    {\(4\)}
    {\(8\)}
\end{problem}

\begin{answer}
C
\end{answer}

\begin{solution}
将抛物线写成标准形式 \(y^2=4px\)，由 \(4p=8\) 得 \(p=2\). 因此焦点为 \((2,0)\)，准线为 \(x=-2\)，焦点到准线的距离为 \(2-(-2)=4\)，故选 C.
\end{solution}

\begin{problem}
一个单位有职工 \(800\text{ 人}\)，其中具有高级职称的 \(160\text{ 人}\)，具有中级职称的 \(320\text{ 人}\)，具有初级职称的 \(200\text{ 人}\)，其余人员 \(120\text{ 人}\). 为了解职工收入情况，决定采用分层抽样的方法，从中抽取容量为 40 的样本. 则从上述各层中依次抽取的人数分别是（\quad）

\choices
    {\(12, 24, 15, 9\)}
    {\(9\)，\(12\)，\(12\)，\(7\)}
    {\(8\)，\(15\)，\(12\)，\(5\)}
    {\(8\)，\(16\)，\(10\)，\(6\)}
\end{problem}

\begin{answer}
D
\end{answer}

\begin{solution}
分层抽样时，各层抽取人数与该层人数成正比. 抽样比例为 \(40/800=1/20\)，所以四层依次抽取 \(160/20=8\)、\(320/20=16\)、\(200/20=10\)、\(120/20=6\) 人，故选 D.
\end{solution}

\begin{problem}
函数 \( f(x) = x^{2} + mx + 1 \) 的图象关于直线 \( x = 1 \) 对称的充要条件是（\quad）

\choices
    {\(m = -2\)}
    {\(m = 2\)}
    {\(m = -1\)}
    {\(m = 1\)}
\end{problem}

\begin{answer}
A
\end{answer}

\begin{solution}
二次函数 \(f(x)=x^2+mx+1\) 的图象是开口向上的抛物线，其对称轴为 \(x=-m/2\). 题设要求对称轴为 \(x=1\)，因此 \(-m/2=1\)，解得 \(m=-2\)，故选 A.
\end{solution}

\begin{problem}
设点 \(M\) 是线段 \(BC\) 的中点，点 \(A\) 在直线 \(BC\) 外，\( \overrightarrow{BC}^{2}=16 \)，\( \left|\overrightarrow{AB}+\overrightarrow{AC}\right|=\left|\overrightarrow{AB}-\overrightarrow{AC}\right| \)，则 \( \overrightarrow{AM}= \)（\quad）

\choices
    {\(8\)}
    {\(4\)}
    {\(2\)}
    {\(1\)}
\end{problem}

\begin{answer}
C
\end{answer}

\begin{solution}
由题设两边平方得
\(|\overrightarrow{AB}|^2+|\overrightarrow{AC}|^2+2\overrightarrow{AB}\cdot\overrightarrow{AC}=|\overrightarrow{AB}|^2+|\overrightarrow{AC}|^2-2\overrightarrow{AB}\cdot\overrightarrow{AC}\)，所以 \(\overrightarrow{AB}\cdot\overrightarrow{AC}=0\). 又因为 \(M\) 是 \(BC\) 的中点，\(\overrightarrow{AM}=\frac12(\overrightarrow{AB}+\overrightarrow{AC})\)，从而
\(|\overrightarrow{AM}|^2=\frac14\left(|\overrightarrow{AB}|^2+|\overrightarrow{AC}|^2\right)\).
另一方面，\(\overrightarrow{BC}=\overrightarrow{AC}-\overrightarrow{AB}\)，故 \(|\overrightarrow{BC}|^2=|\overrightarrow{AC}|^2+|\overrightarrow{AB}|^2=16\). 因此 \(|\overrightarrow{AM}|=\frac12\sqrt{16}=2\)，故选 C.
\end{solution}

\begin{problem}
将函数 \( y = \sin x \) 的图象上所有的点向右平行移动 \( \frac{\pi}{10} \) 个单位长度，再把所得各点的横坐标伸长到原来的 2 倍（纵坐标不变），所得图象的函数解析式是（\quad）

\choices
    {\( y = \sin \left(2x - \frac{\pi}{10}\right) \)}
    {\( y = \sin \left(2x - \frac{\pi}{5}\right) \)}
    {\( y = \sin \left(\frac{1}{2}x - \frac{\pi}{10}\right) \)}
    {\( y = \sin \left(\frac{1}{2}x - \frac{\pi}{20}\right) \)}
\end{problem}

\begin{answer}
C
\end{answer}

\begin{solution}
图象向右平移 \(\pi/10\) 后，函数变为 \(y=\sin\left(x-\frac{\pi}{10}\right)\). 横坐标伸长到原来的 \(2\) 倍，相当于把新函数中的 \(x\) 替换为 \(x/2\)，所以所得函数为 \(y=\sin\left(\frac{x}{2}-\frac{\pi}{10}\right)\)，故选 C.
\end{solution}

\begin{problem}
某加工厂用某原料由甲车间加工出 A 产品，由乙车间加工出 B 产品. 甲车间加工一箱原料需耗费工时 \(10\text{ 小时}\) 可加工出 \(7\text{ 千克}\) A 产品，每千克 A 产品获利 \(40\text{ 元}\). 乙车间加工一箱原料需耗费工时 \(6\text{ 小时}\) 可加工出 \(4\text{ 千克}\) B 产品，每千克 B 产品获利 \(50\text{ 元}\). 甲，乙两车间每天共能完成至多 \(70\text{ 箱}\) 原料的加工. 每天甲，乙两车间耗费工时总和不得超过 \(480\text{ 小时}\)，甲，乙两车间每天总获利最大的生产计划为（\quad）

\choices
    {甲车间加工原料 \(10\) 箱，乙车间加工原料 \(60\) 箱}
    {甲车间加工原料 \(15\) 箱，乙车间加工原料 \(55\) 箱}
    {甲车间加工原料 \(18\) 箱，乙车间加工原料 \(50\) 箱}
    {甲车间加工原料 \(40\) 箱，乙车间加工原料 \(30\) 箱}
\end{problem}

\begin{answer}
B
\end{answer}

\begin{solution}
设甲、乙两车间每天分别加工 \(x\)、\(y\) 箱原料，则 \(x\geq0\)、\(y\geq0\)，且 \(x+y\leq70\)、\(10x+6y\leq480\). 每加工一箱甲车间原料获利 \(7\times40=280\) 元，每加工一箱乙车间原料获利 \(4\times50=200\) 元，因此总获利为 \(P=280x+200y\).

可行域的顶点为 \((0,0)\)、\((0,70)\)、\((15,55)\) 和 \((48,0)\). 在直线 \(x+y=70\) 上，\(P=280x+200(70-x)=80x+14000\)，随 \(x\) 增大而增大；工时约束与该直线的交点由 \(10x+6(70-x)=480\) 得为 \((15,55)\). 逐点计算得
\(P(0,0)=0\)、\(P(0,70)=14000\)、\(P(15,55)=15200\)、\(P(48,0)=13440\). 所以最大获利在 \((15,55)\) 处取得，最优生产计划为甲加工 \(15\) 箱、乙加工 \(55\) 箱，故选 B.
\end{solution}

\begin{problem}
由 1，2，3，4，5 组成没有重复数字且 1，2 都不与 5 相邻的五位数的个数是（\quad）

\choices
    {\(36\)}
    {\(32\)}
    {\(28\)}
    {\(24\)}
\end{problem}

\begin{answer}
A
\end{answer}

\begin{solution}
先从全部排列中减去不合条件的排列. 设 \(E_1\) 表示数字 \(1\) 与 \(5\) 相邻，\(E_2\) 表示数字 \(2\) 与 \(5\) 相邻. 总排列数为 \(5!=120\)，而 \(|E_1|=|E_2|=2\times4!=48\).

在 \(E_1\cap E_2\) 中，数字 \(5\) 可以位于第 \(2\)、第 \(3\) 或第 \(4\) 位. 对每一种位置，\(1\) 和 \(2\) 分别排在 \(5\) 两侧有 \(2\) 种，剩余两个数字有 \(2\) 种，共 \(3\times2\times2=12\) 种. 因此满足条件的五位数有 \(120-48-48+12=36\) 个，故选 A.
\end{solution}

\begin{problem}
椭圆 \(\frac{x^2}{a^2} + \frac{y^2}{b^2} = 1\)（\(a > b > 0\)）的右焦点为 \(F\)，其右准线与 \(x\) 轴的交点为 \(A\)，在椭圆上存在点 \(P\) 满足线段 \(AP\) 的垂直平分线过点 \(F\)，则椭圆离心率的取值范围是（\quad）

\choices
    {\(\left(0, \frac{\sqrt{2}}{2}\right]\)}
    {\(\left(0, \frac{1}{2}\right]\)}
    {\(\left[\sqrt{2} - 1, 1\right)\)}
    {\(\left[\frac{1}{2}, 1\right)\)}
\end{problem}

\begin{answer}
D
\end{answer}

\begin{solution}
设椭圆焦距参数为 \(c=\sqrt{a^2-b^2}\)，离心率为 \(e=c/a\)，则 \(0<e<1\). 右焦点和右准线与 \(x\) 轴的交点分别为 \(F(c,0)\) 和 \(A(a^2/c,0)\). 因为 \(F\) 在线段 \(AP\) 的垂直平分线上，所以 \(FA=FP\). 计算得
\(FA=\frac{a^2}{c}-c=\frac{a^2-c^2}{c}=\frac{b^2}{c}=\frac{a(1-e^2)}{e}\).

对椭圆上的点 \(P\)，距离 \(FP\) 的取值范围为 \([a-c,a+c]=[a(1-e),a(1+e)]\). 因此存在这样的 \(P\) 的充要条件是 \(a(1-e)\leq\frac{a(1-e^2)}{e}\leq a(1+e)\). 左侧不等式对 \(0<e<1\) 恒成立；右侧不等式化为 \(1-e\leq e\)，即 \(e\geq1/2\). 所以离心率的取值范围为 \([1/2,1)\)，故选 D.
\end{solution}

\begin{problem}
设 \(a > b > 0\)，则 \(a^2 + \frac{1}{ab} + \frac{1}{a(a - b)}\) 的最小值是（\quad）

\choices
    {\(1\)}
    {\(2\)}
    {\(3\)}
    {\(4\)}
\end{problem}

\begin{answer}
D
\end{answer}

\begin{solution}
令 \(x=b/a\)，则 \(0<x<1\)，并令 \(u=a^2>0\). 原式化为 \(u+\frac{1}{ux}+\frac{1}{u(1-x)}=u+\frac{1}{ux(1-x)}\). 由于 \(x(1-x)\leq1/4\)，有 \(1/[x(1-x)]\geq4\). 由基本不等式，
\(u+\frac{1}{ux(1-x)}\geq2\sqrt{\frac{1}{x(1-x)}}\geq4\).

当 \(x=1/2\) 且 \(u=2\) 时，即 \(b/a=1/2\)、\(a^2=2\)，上式两次取等号，满足 \(a>b>0\). 因此原式的最小值为 \(4\)，故选 D.
\end{solution}

\begin{problem}
半径为 \(R\) 的球 \(O\) 的直径 \(AB\) 垂直于平面 \(\alpha\)，垂足为 \(B\)，\(\triangle BCD\) 是平面 \(\alpha\) 内边长为 \(R\) 的正三角形，线段 \(AC\)，\(AD\) 分别与球面交于点 \(M\)，\(N\). 那么 \(M\)，\(N\) 两点间的球面距离是（\quad）

\bitmapfigure[max width=0.38\linewidth]{img/2010/sichuan_liberal/q12.png}

\choices
    {\(R\arccos\frac{17}{25}\)}
    {\(R\arccos\frac{18}{25}\)}
    {\(\frac{4}{3}\pi R\)}
    {\(\frac{4}{15}\pi R\)}
\end{problem}

\begin{answer}
A
\end{answer}

\begin{solution}
取 \(O\) 为原点，令 \(A=(0,0,R)\)、\(B=(0,0,-R)\)，从而平面 \(\alpha\) 可取为 \(z=-R\). 由正三角形的对称性，不妨设 \(C=(R,0,-R)\)、\(D=\left(\frac R2,\frac{\sqrt3R}{2},-R\right)\).

直线 \(AC\) 上的点可写为 \(A+s(C-A)=(sR,0,R-2sR)\). 代入球面方程得 \(s^2R^2+(R-2sR)^2=R^2\)，即 \(s(5s-4)=0\). 除去 \(s=0\) 对应的点 \(A\)，得 \(s=4/5\)，所以 \(M=\left(\frac{4R}{5},0,-\frac{3R}{5}\right)\).

直线 \(AD\) 上的点可写为 \(A+s(D-A)=\left(\frac{sR}{2},\frac{\sqrt3sR}{2},R-2sR\right)\). 代入球面方程得
\(\frac{s^2R^2}{4}+\frac{3s^2R^2}{4}+(R-2sR)^2=R^2\)，即 \(s(5s-4)=0\). 除去 \(s=0\) 对应的点 \(A\)，得 \(s=4/5\)，所以 \(N=\left(\frac{2R}{5},\frac{2\sqrt3R}{5},-\frac{3R}{5}\right)\).

于是 \(\overrightarrow{OM}\cdot\overrightarrow{ON}=\frac{8R^2}{25}+\frac{9R^2}{25}=\frac{17R^2}{25}\)，且 \(|OM|=|ON|=R\). 设小圆弧 \(MN\) 所对的圆心角为 \(\theta\)，则 \(\cos\theta=\frac{\overrightarrow{OM}\cdot\overrightarrow{ON}}{|OM||ON|}=\frac{17}{25}\). 因此球面距离为 \(R\theta=R\arccos\frac{17}{25}\)，故选 A.
\end{solution}

\section{填空题}

\begin{problem}
\(\left(x - \frac{2}{x}\right)^4\) 的展开式中的常数项为\fillinblank{}. （用数字作答）

\end{problem}

\begin{answer}
24
\end{answer}

\begin{solution}
展开式的一般项为 \(\mathrm C_4^k x^{4-k}\left(-\frac2x\right)^k=\mathrm C_4^k(-2)^k x^{4-2k}\). 常数项要求 \(4-2k=0\)，所以 \(k=2\)，常数项为 \(\mathrm C_4^2(-2)^2=6\times4=24\).
\end{solution}

\begin{problem}
直线 \(x - 2y + 5 = 0\) 与圆 \(x^{2} + y^{2} = 8\) 相交于 \(A\)，\(B\) 两点，则 \(|AB| = \)\fillinblank{}.
\end{problem}

\begin{answer}
\(2\sqrt{3}\)
\end{answer}

\begin{solution}
圆心为原点，半径为 \(\sqrt8=2\sqrt2\). 原点到直线 \(x-2y+5=0\) 的距离为 \(d=\frac5{\sqrt{1^2+(-2)^2}}=\sqrt5\). 设弦 \(AB\) 的中点为 \(H\)，则 \(OH\perp AB\)，并且 \(AH=\sqrt{(2\sqrt2)^2-(\sqrt5)^2}=\sqrt3\). 所以 \(|AB|=2AH=2\sqrt3\).
\end{solution}

\begin{problem}
如图，二面角 \(\alpha - l - \beta\) 的大小是 \(60^{\circ}\)，\(AB \subset \alpha\)，\(B \in l\)，\(AB\) 与 \(l\) 所成的角为 \(30^{\circ}\)，则 \(AB\) 与平面 \(\beta\) 所成角的正弦值是\fillinblank{}.

\bitmapfigure[max width=0.46\linewidth]{img/2010/sichuan_liberal/q15.png}
\end{problem}

\begin{answer}
\(\frac{\sqrt{3}}{4}\)
\end{answer}

\begin{solution}
取 \(l\) 为 \(x\) 轴，令平面 \(\alpha\) 为 \(z=0\). 二面角为 \(60^\circ\)，所以平面 \(\beta\) 可看作平面 \(\alpha\) 绕 \(l\) 旋转 \(60^\circ\) 所得的平面. 取 \(AB\) 的单位方向向量为 \(\bs{u}=(\cos30^\circ,\sin30^\circ,0)\)，取平面 \(\beta\) 的单位法向量为 \(\bs{n}=(0,\sin60^\circ,\cos60^\circ)\).

设 \(AB\) 与平面 \(\beta\) 所成的角为 \(\theta\)，则 \(\sin\theta=|\bs{u}\cdot\bs{n}|\). 因而 \(\sin\theta=\sin30^\circ\sin60^\circ=\frac12\times\frac{\sqrt3}{2}=\frac{\sqrt3}{4}\).
\end{solution}

\begin{problem}
设 \(S\) 为实数集 \(\R\) 的非空子集，若对任意 \(x\)，\(y \in S\)，都有 \(x + y\)，\(x - y\)，\(xy \in S\)，则称 \(S\) 为封闭集. 下列命题：

\begin{circlelist}
\item 集合 \(S = \{a + b\sqrt{3} \mid a, b\text{ 为整数}\}\) 为封闭集；
\item 若 \(S\) 为封闭集，则一定有 \(0 \in S\)；
\item 封闭集一定是无限集；
\item 若 \(S\) 为封闭集，则满足 \(S \subseteq T \subseteq \R\) 的任意集合 \(T\) 也是封闭集.
\end{circlelist}

其中的真命题是\fillinblank{}（写出所有真命题的序号）
\end{problem}

\begin{answer}
\(\circled{1}\circled{2}\)
\end{answer}

\begin{solution}
对于命题 \(\circled{1}\)，任取 \(a+b\sqrt3\) 与 \(c+d\sqrt3\)，其中 \(a\)，\(b\)，\(c\)，\(d\) 为整数，则
\((a+b\sqrt3)+(c+d\sqrt3)=(a+c)+(b+d)\sqrt3\)，
\((a+b\sqrt3)-(c+d\sqrt3)=(a-c)+(b-d)\sqrt3\)，
以及 \((a+b\sqrt3)(c+d\sqrt3)=(ac+3bd)+(ad+bc)\sqrt3\). 三式仍属于该集合，所以命题 \(\circled{1}\) 正确.

对于命题 \(\circled{2}\)，封闭集非空，任取 \(x\in S\)，由差的封闭性得 \(x-x=0\in S\)，所以命题 \(\circled{2}\) 正确.

命题 \(\circled{3}\) 不正确，因为 \(S=\{0\}\) 非空，且对其中任意两个元素进行加、减、乘运算所得仍为 \(0\)，所以它是有限封闭集. 命题 \(\circled{4}\) 也不正确，例如取封闭集 \(S=\{0\}\)，取 \(T=\{0,1\}\)，则 \(S\subseteq T\subseteq\R\)，但 \(1+1=2\notin T\)，所以 \(T\) 不是封闭集. 因此所有真命题的序号为 \(\circled{1}\circled{2}\).
\end{solution}

\section{解答题}

\begin{problem}
某种有奖销售的饮料，瓶盖内印有“奖励一瓶”或“谢谢购买”字样，购买一瓶若其瓶盖内印有“奖励一瓶”字样即为中奖，中奖概率为 \(\frac{1}{6}\). 甲，乙，丙三位同学每人购买了一瓶该饮料.

\begin{enumerate}
\item 求三位同学都没有中奖的概率；
\item 求三位同学中至少有两位没有中奖的概率.
\end{enumerate}
\end{problem}

\begin{solution}
\begin{enumerate}
\item 各瓶盖的抽奖结果相互独立. 每位同学没有中奖的概率为 \(1-\frac16=\frac56\)，三位同学都没有中奖的概率为 \(\left(\frac56\right)^3=\frac{125}{216}\).
\item 三位同学中至少有两位没有中奖包括恰有两位没有中奖和三位都没有中奖两种互斥情况，所以所求概率为
\( \mathrm C_3^2\left(\frac56\right)^2\frac16+\left(\frac56\right)^3=\frac{75}{216}+\frac{125}{216}=\frac{25}{27}\).
\end{enumerate}
\end{solution}

\begin{problem}
已知正方体 \(ABCD - A'B'C'D'\) 的棱长为 1，点 \(M\) 是棱 \(AA'\) 的中点，点 \(O\) 是对角线 \(BD'\) 的中点.

\begin{enumerate}
\item 求证：\(OM\) 为异面直线 \(AA'\) 和 \(BD'\) 的公垂线；
\item 求二面角 \(M - BC' - B'\) 的大小.
\end{enumerate}

\bitmapfigure[max width=0.38\linewidth]{img/2010/sichuan_liberal/q18.png}
\end{problem}

\begin{solution}
\begin{enumerate}
\item 建立直角坐标系，令 \(A=(0,0,0)\)、\(B=(1,0,0)\)、\(C=(1,1,0)\)、\(D=(0,1,0)\)，并令 \(A'\)、\(B'\)、\(C'\)、\(D'\) 的竖直坐标都为 \(1\). 则 \(M=(0,0,\frac12)\)，\(D'=(0,1,1)\)，所以 \(O\)，即 \(BD'\) 的中点为 \(O=(\frac12,\frac12,\frac12)\).

于是 \(\overrightarrow{MO}=(\frac12,\frac12,0)\)，而 \(AA'\) 的方向向量为 \((0,0,1)\)，故 \(\overrightarrow{MO}\cdot(0,0,1)=0\)，从而 \(OM\perp AA'\). 又 \(BD'\) 的方向向量为 \((-1,1,1)\)，有 \(\overrightarrow{MO}\cdot(-1,1,1)=0\)，从而 \(OM\perp BD'\). 因为 \(M\in AA'\)、\(O\in BD'\)，所以 \(OM\) 是两异面直线 \(AA'\) 和 \(BD'\) 的公垂线.
\item 平面 \(B'C'B\) 的两个方向向量可取 \(\overrightarrow{BC'}=(0,1,1)\)、\(\overrightarrow{BB'}=(0,0,1)\)，故其法向量可取 \(\bs{n}_2=(1,0,0)\). 平面 \(MBC'\) 的两个方向向量可取 \(\overrightarrow{BC'}=(0,1,1)\)、\(\overrightarrow{BM}=(-1,0,\frac12)\)，故其法向量可取
\(\bs{n}_1=\overrightarrow{BC'}\times\overrightarrow{BM}=(\frac12,-1,1)\).

设二面角的平面角为 \(\theta\)，则 \(\theta\) 是两法向量的锐角，因而
\(\cos\theta=\frac{|\bs{n}_1\cdot\bs{n}_2|}{|\bs{n}_1||\bs{n}_2|}=\frac{1/2}{3/2}=\frac13\). 所以 \(\tan\theta=\sqrt{1-\frac19}/(\frac13)=2\sqrt2\)，即 \(\theta=\arctan(2\sqrt2)\).
\end{enumerate}
\end{solution}

\begin{problem}
\begin{enumerate}
\item
\begin{enumerate}
\item 证明两角和的余弦公式 \(C_{(\alpha+\beta)}\)：\(\cos (\alpha + \beta) = \cos \alpha \cos \beta - \sin \alpha \sin \beta\)；
\item 由 \(C_{(\alpha+\beta)}\) 推导两角和的正弦公式 \(S_{(\alpha+\beta)}\)：\(\sin (\alpha + \beta) = \sin \alpha \cos \beta + \cos \alpha \sin \beta\).
\end{enumerate}
\item 已知 \(\cos \alpha = -\frac{4}{5}\)，\(\alpha \in \left(\pi, \frac{3\pi}{2}\right)\)，\(\tan \beta = -\frac{1}{3}\)，\(\beta \in \left(\frac{\pi}{2}, \pi\right)\)，求 \(\cos (\alpha + \beta)\).
\end{enumerate}
\end{problem}

\begin{solution}
\begin{enumerate}
\item
\begin{enumerate}
\item 在单位圆中，令角 \(\alpha\) 对应的点为 \(P=(\cos\alpha,\sin\alpha)\). 将点 \(P\) 绕原点逆时针旋转 \(\beta\)，所得点对应角 \(\alpha+\beta\). 旋转后，原来的单位向量 \((1,0)\) 变为 \((\cos\beta,\sin\beta)\)，原来的单位向量 \((0,1)\) 变为 \((-\sin\beta,\cos\beta)\). 因此旋转后的点坐标为
\(\cos\alpha(\cos\beta,\sin\beta)+\sin\alpha(-\sin\beta,\cos\beta)\)
\(=(\cos\alpha\cos\beta-\sin\alpha\sin\beta,\ \cos\alpha\sin\beta+\sin\alpha\cos\beta)\).
另一方面，该点坐标也为 \((\cos(\alpha+\beta),\sin(\alpha+\beta))\)，比较第一坐标得 \(\cos(\alpha+\beta)=\cos\alpha\cos\beta-\sin\alpha\sin\beta\).
\item 由余弦公式 \(C_{(\alpha+\beta)}\) 及诱导关系 \(\sin\theta=\cos\left(\frac{\pi}{2}-\theta\right)\)，有 \(\sin(\alpha+\beta)=\cos\left(\frac{\pi}{2}-(\alpha+\beta)\right)=\cos\left(\left(\frac{\pi}{2}-\alpha\right)+(-\beta)\right)\). 应用余弦公式，并利用 \(\cos\left(\frac{\pi}{2}-\alpha\right)=\sin\alpha\)、\(\sin\left(\frac{\pi}{2}-\alpha\right)=\cos\alpha\)、\(\cos(-\beta)=\cos\beta\)、\(\sin(-\beta)=-\sin\beta\)，得 \(\sin(\alpha+\beta)=\sin\alpha\cos\beta+\cos\alpha\sin\beta\).
\end{enumerate}
\item 因为 \(\alpha\in(\pi,\frac{3\pi}{2})\)，所以 \(\sin\alpha<0\)，从而 \(\sin\alpha=-\sqrt{1-\cos^2\alpha}=-\frac35\). 又因为 \(\beta\in(\frac\pi2,\pi)\)，所以 \(\sin\beta>0\)、\(\cos\beta<0\). 由 \(\tan\beta=-1/3\) 得
\(\sin\beta=\frac1{\sqrt{10}}=\frac{\sqrt{10}}{10}\)，\(\cos\beta=-\frac3{\sqrt{10}}=-\frac{3\sqrt{10}}{10}\).

代入两角和的余弦公式，得
\(\cos(\alpha+\beta)=\left(-\frac45\right)\left(-\frac3{\sqrt{10}}\right)-\left(-\frac35\right)\left(\frac1{\sqrt{10}}\right)=\frac3{\sqrt{10}}=\frac{3\sqrt{10}}{10}\).
\end{enumerate}
\end{solution}

\begin{problem}
已知等差数列 \(\{a_n\}\) 的前3项和为6，前8项和为\(-4\).

\begin{enumerate}
\item 求数列 \( \{a_n\} \) 的通项公式；
\item 设 \( b_{n} = (4 - a_{n})q^{n - 1} \) （\(q \neq 0\)，\(n \in \N^*\)），求数列 \( \{b_{n}\} \) 的前 \( n \) 项和 \( S_{n} \).
\end{enumerate}
\end{problem}

\begin{solution}
\begin{enumerate}
\item 设等差数列首项为 \(a_1=a\)，公差为 \(d\). 由前 \(3\) 项和为 \(6\)，得 \(3(a+d)=6\)，即 \(a+d=2\). 由前 \(8\) 项和为 \(-4\)，得 \(4(2a+7d)=-4\)，即 \(2a+7d=-1\). 联立解得 \(d=-1\)、\(a=3\)，所以 \(a_n=a+(n-1)d=4-n\).
\item 由 \(a_n=4-n\) 得 \(4-a_n=n\)，所以 \(b_n=nq^{n-1}\). 当 \(q\ne1\) 时，令 \(S_n=\sum_{k=1}^n kq^{k-1}\)，则
\((1-q)S_n=1+q+q^2+\cdots+q^{n-1}-nq^n\).
再利用 \(1+q+\cdots+q^{n-1}=(1-q^n)/(1-q)\)，得
\(S_n=\frac{1-(n+1)q^n+nq^{n+1}}{(1-q)^2}=\frac{nq^{n+1}-(n+1)q^n+1}{(q-1)^2}\).
当 \(q=1\) 时，\(S_n=1+2+\cdots+n=\frac{n(n+1)}2\).
\end{enumerate}
\end{solution}

\begin{problem}
已知定点 \(A(-1,0)\)，\(F(2,0)\). 定直线 \(l: x = \frac{1}{2}\)，不在 \(x\) 轴上的动点 \(P\) 与点 \(F\) 的距离是它到直线 \(l\) 的距离的 2 倍. 设点 \(P\) 的轨迹为 \(E\)；过点 \(F\) 的直线交 \(E\) 于 \(B\)，\(C\) 两点，直线 \(AB\)，\(AC\) 分别交 \(l\) 于点 \(M\)，\(N\).

\begin{enumerate}
\item 求 \(E\) 的方程；
\item 试判断以线段 \(MN\) 为直径的圆是否过点 \(F\)，并说明理由.
\end{enumerate}
\end{problem}

\begin{solution}
\begin{enumerate}
\item 设 \(P=(x,y)\). 由题意，\(P\) 到 \(F(2,0)\) 的距离是它到直线 \(l:x=\frac12\) 的距离的 \(2\) 倍，所以
\((x-2)^2+y^2=4\left(x-\frac12\right)^2\).
化简得 \(3x^2-y^2=3\)，即 \(x^2-\frac{y^2}{3}=1\). 结合题设 \(P\) 不在 \(x\) 轴上，轨迹为 \(E: x^2-\frac{y^2}{3}=1\)，并满足 \(y\ne0\).
\item 设过 \(F\) 的直线不是竖直直线，写成 \(y=k(x-2)\). 因为 \(B\)、\(C\) 不在 \(x\) 轴上，所以 \(k\ne0\). 设它与 \(E\) 的交点 \(B\)、\(C\) 的横坐标分别为 \(x_1\)、\(x_2\)，则 \(x_1\)、\(x_2\) 是方程
\((3-k^2)x^2+4k^2x-(4k^2+3)=0\)
的两个根. 由于 \(B\)、\(C\) 是两个交点，必须有 \(3-k^2\ne0\)，由韦达定理
\(x_1+x_2=-\frac{4k^2}{3-k^2}\)，\(x_1x_2=-\frac{4k^2+3}{3-k^2}\).

直线 \(A(-1,0)B\) 与 \(x=\frac12\) 的交点记为 \(M\)，直线 \(A C\) 与该直线的交点记为 \(N\). 由相似三角形或直线参数方程，二者的纵坐标分别为
\(y_M=\frac{3k(x_1-2)}{2(x_1+1)}\)、\(y_N=\frac{3k(x_2-2)}{2(x_2+1)}\).
由根与系数的关系，
\(\frac{(x_1-2)(x_2-2)}{(x_1+1)(x_2+1)}
=\frac{x_1x_2-2(x_1+x_2)+4}{x_1x_2+(x_1+x_2)+1}=-\frac1{k^2}\).
所以 \(y_My_N=-\frac94\). 因为 \(M=(\frac12,y_M)\)、\(N=(\frac12,y_N)\)、\(F=(2,0)\)，有
\(\overrightarrow{FM}\cdot\overrightarrow{FN}=\left(-\frac32,y_M\right)\cdot\left(-\frac32,y_N\right)=\frac94+y_My_N=0\).
故 \(\angle MFN=90^\circ\)，点 \(F\) 在以 \(MN\) 为直径的圆上.

若过 \(F\) 的直线为竖直直线 \(x=2\)，则其与 \(E\) 的交点为 \((2,3)\)、\((2,-3)\)，相应地 \(M=(\frac12,\frac32)\)、\(N=(\frac12,-\frac32)\)，同样有 \(\overrightarrow{FM}\cdot\overrightarrow{FN}=0\). 因此无论何种过 \(F\) 且与 \(E\) 有两个交点的直线，所述圆都经过 \(F\).
\end{enumerate}
\end{solution}

\begin{problem}
设 \( f(x) = \frac{1 + a^x}{1 - a^x} \) （\(a > 0\) 且 \(a \neq 1\)），\( g(x) \) 是 \( f(x) \) 的反函数.

\begin{enumerate}
\item 求 \( g(x) \).
\item 当 \(x \in [2,6]\) 时，恒有 \(g(x) > \log_a \frac{t}{(x^2 - 1)(7 - x)}\) 成立，求 \(t\) 的取值范围；
\item 当 \(0 < a \leqslant \frac{1}{2}\) 时，试比较 \(f(1) + f(2) + \cdots + f(n)\) 与 \(n + 4\) 的大小，并说明理由.
\end{enumerate}
\end{problem}

\begin{solution}
\begin{enumerate}
\item 设 \(y=f(x)\)，则 \(y(1-a^x)=1+a^x\)，所以 \(a^x=(y-1)/(y+1)\). 因为 \(a^x>0\) 且 \(a^x\ne1\)，所以 \(y\in(-\infty,-1)\cup(1,+\infty)\)，从而
\(g(x)=\log_a\frac{x-1}{x+1}\)，定义域为 \((-\infty,-1)\cup(1,+\infty)\).
\item 令 \(\varphi(x)=(x^2-1)(7-x)\frac{x-1}{x+1}=(x-1)^2(7-x)\). 对 \(x\in[2,6]\)，有 \((x^2-1)(7-x)>0\).

当 \(a>1\) 时，对数函数单调递增，题设不等式等价于
\(\frac{x-1}{x+1}>\frac{t}{(x^2-1)(7-x)}\)，即 \(t<\varphi(x)\) 对一切 \(x\in[2,6]\) 成立. 又 \(\varphi'(x)=3(x-1)(5-x)\)，所以 \(\varphi\) 在 \([2,5]\) 上递增，在 \([5,6]\) 上递减，且 \(\varphi(2)=5\)、\(\varphi(5)=32\)、\(\varphi(6)=25\). 因而 \(0<t<5\).

当 \(0<a<1\) 时，对数函数单调递减，题设不等式等价于 \(t>\varphi(x)\) 对一切 \(x\in[2,6]\) 成立. 因此 \(t>\max\varphi(x)=32\)，即 \(t>32\).
\item 当 \(0<a\leq1/2\) 时，对 \(k\geq1\) 有
\(f(k)=1+\frac{2a^k}{1-a^k}\leq1+\frac{2(1/2)^k}{1-(1/2)^k}=1+\frac{2}{2^k-1}\).
于是
\(\sum_{k=1}^n f(k)\leq n+2+\sum_{k=2}^n\frac{2}{2^k-1}\).
对 \(k\geq2\)，有 \(2^k-1>2^{k-1}\)，所以 \(\frac{2}{2^k-1}<\frac{4}{2^k}\). 因而
\(\sum_{k=1}^n f(k)<n+2+\sum_{k=2}^{\infty}\frac{4}{2^k}=n+2+2=n+4\).
所以 \(f(1)+f(2)+\cdots+f(n)<n+4\).
\end{enumerate}
\end{solution}
